\section{Conventions and the difference lemma}\label{sec:conv}

These conventions hold throughout the paper.
\begin{enumerate}[(1),leftmargin=*]
\item An $h\times m$ grid has $h$ rows and $m$ columns, so $hm$ points. A $2\times m$
grid is a \emph{strip}, and an $n\times n$ grid is a \emph{square}.
\item The point $(i,j)$ is in column $i$ and row $j$, and $(1,1)$ is the top-left
point. So $i$ grows to the right and $j$ grows downward, as in
Figure~\ref{fig:game}(a). The set of all points is $V$.
\item A \emph{rectangle} is a rectangle of positive area whose four corners are grid
points, tilted ones included. We identify $R$ with its corner set $C(R)$. A
rectangle is \emph{axis-aligned} if its sides are parallel to the rows and columns,
and \emph{tilted} otherwise. Four distinct points form a rectangle if and only if
they can be paired into two segments with the same midpoint and the same length.
\item A \emph{blackout} is a set $S\subseteq V$, and the \emph{kept set} is
$K=V\setminus S$. The set $B_i$ of \emph{black rows} of
column $i$ consists of the rows $j$ with $(i,j)\in S$. A column is \emph{empty}, a
\emph{singleton}, a \emph{double}, or a \emph{triple} when $|B_i|$ is $0$, $1$, $2$,
or $3$, and it is \emph{full} when $|B_i|=h$. The \emph{pair set} $F$ is the set of
row pairs $\{r,s\}$ with $\{r,s\}\subseteq B_i$ for some column $i$. We put
$q_h=\lfloor h^2/4\rfloor$.
\item For a blackout $S$ and a rectangle $R$, the \emph{presentation} of $R$ is the
set of its white corners, $\pres(R)=C(R)\setminus S$.
\end{enumerate}

\begin{definition}[valid blackout]\label{def:valid}
A blackout $S$ is \emph{valid} if different rectangles always have different
presentations, that is, if $\pres(R_1)\ne\pres(R_2)$ whenever $R_1\ne R_2$. The
\emph{maximum blackout} $f(h,m)$ of the $h\times m$ grid is the largest size of a
valid blackout, and $a_h(m)$ is the number of valid blackouts of size $f(h,m)$.
Blackouts are counted with their coordinates, not up to the symmetries of the grid.
\end{definition}

Swapping rows and columns maps rectangles to rectangles, so $f(h,m)=f(m,h)$ and
$a_h(m)=a_m(h)$. In the OEIS, $f(h,m)$ is \seqnum{A397315} and $a_h(m)$ is
\seqnum{A399580}, both read by antidiagonals with $h,m\ge2$. Section~\ref{sec:height}
also uses width $0$ and height $1$, where there are no rectangles, so $f(h,0)=0$,
$a_h(0)=1$, $f(1,m)=m$, and $a_1(m)=1$.

\begin{remark}[the empty presentation]\label{rem:empty}
The empty presentation is allowed. It means that every corner of the rectangle is
black. Two hidden rectangles would share it, so a valid blackout hides at most one
rectangle. This is a consequence of the definition, not an extra rule.
\end{remark}

The next lemma says when two rectangles look the same. Here
$A\triangle B=(A\setminus B)\cup(B\setminus A)$ is the symmetric difference.

\begin{lemma}[difference sets]\label{lem:diff}
For all sets $A$, $B$, and $S$ we have $A\setminus S=B\setminus S$ if and only if
$A\,\triangle\,B\subseteq S$.
\end{lemma}

\begin{proof}
A point lies in $(A\setminus S)\triangle(B\setminus S)$ exactly when it is outside
$S$ and in exactly one of $A$ and $B$. So
$(A\setminus S)\triangle(B\setminus S)=(A\triangle B)\setminus S$. Two sets are
equal exactly when their symmetric difference is empty, and
$(A\triangle B)\setminus S$ is empty exactly when $A\triangle B\subseteq S$.
\end{proof}

\begin{corollary}[hitting sets]\label{cor:hitting}
A blackout $S$ is valid if and only if the kept set $K=V\setminus S$ contains at
least one point of $C(R_1)\triangle C(R_2)$ for every pair of distinct rectangles
$R_1$, $R_2$. So $f(h,m)$ is $hm$ minus the smallest size of a set of points that
meets every such difference set.
\end{corollary}

\begin{proof}
Apply Lemma~\ref{lem:diff} to $C(R_1)$, $C(R_2)$, and $S$ for every pair.
\end{proof}

A difference set that contains another one is met whenever the smaller one is met.
So only the minimal difference sets matter, and the programs of
Section~\ref{sec:verify} delete the others before they search.

\begin{lemma}[three corners suffice]\label{lem:threecorners}
Three corners of a rectangle determine the rectangle. So if two rectangles have the
same presentation and the presentation has at least three points, the two rectangles
are equal.
\end{lemma}

\begin{proof}
Let $p$, $q$, and $r$ be three corners of a rectangle. Exactly one of them, say $q$,
is adjacent to both others, and the triangle $pqr$ has its right angle at $q$. A
triangle has at most one right angle, so the three points determine $q$, and the
fourth corner is $p+r-q$. A presentation with at least three points contains three
corners of each rectangle, the same three points, so the two rectangles have the
same corners.
\end{proof}

So two distinct rectangles can only be confused if both show at most two corners.
Section~\ref{sec:graph} gives such rectangles a name.
